\documentclass{article}
\usepackage[T1]{fontenc}
\usepackage{lmodern}
\usepackage{graphicx}
\usepackage{amsmath,amssymb}
\usepackage[a4paper,margin=2cm]{geometry}
\usepackage[absolute,overlay]{textpos}
\setlength{\TPHorizModule}{1mm}
\setlength{\TPVertModule}{1mm}
\usepackage[indonesian]{babel}
\usepackage{hyperref}
\setlength{\emergencystretch}{2em}
% unit: Halaman 21

\begin{document}

% === Halaman 21 (python/native) ===

Jelaslah Kurva Eliptik membentuk Grup <G, +>

Karena:

• Himpunan G: semua titik P(x,y) pada kurva eliptik

• Operasi biner: +

• Semua aksioma terpenuhi sbb:

 1. Closure: semua operasi P + Q berada di dalam G

 2. Asosiatif:  P + (Q + R) = (P + Q) + R

 3. Elemen netral adalah O:   P + O = O + P = P

 4. Elemen invers adalah -P:  P + (-P) = O

 5. Komutatif: P + Q = Q + P     (abelian)

21

\end{document}
